% ============================================================
%  01 封面
% ============================================================
\begin{frame}[plain]
\begin{canvas}
  % 右侧：人群图装饰
  \begin{scope}[shift={(7.6,0.2)}]
    \foreach \n/\x/\y in {a/1.0/5.2, b/2.3/6.1, c/3.6/5.4, d/5.2/6.0, e/6.3/4.9,
                          f/0.6/3.6, g/2.0/4.2, h/3.4/3.7, i/4.9/4.2, j/6.0/3.1,
                          k/1.2/1.9, l/2.7/2.4, m/4.1/1.9, n/5.4/1.4, o/3.0/0.8, p/6.3/1.9}
      \coordinate (\n) at (\x,\y);
    \foreach \u/\v/\w/\o in {a/b/1.2/.55, b/c/.6/.35, c/d/1.4/.6, d/e/.7/.4, a/g/.9/.45, b/g/.5/.3,
                             c/h/1.1/.5, c/i/.6/.35, d/i/1.0/.5, e/i/.5/.3, e/j/1.3/.55, f/g/.8/.4,
                             g/h/1.5/.6, h/i/.7/.4, i/j/.9/.45, f/k/1.0/.5, g/l/.6/.35, h/l/1.1/.5,
                             h/m/.6/.3, i/m/.8/.4, j/n/1.0/.5, j/p/.6/.35, k/l/.9/.45, l/m/1.3/.55,
                             m/n/.7/.4, l/o/.6/.35, m/o/1.0/.5, n/p/1.1/.5, k/o/.5/.3}
      \draw[Navy,opacity=\o,line width=\w pt] (\u) -- (\v);
    \foreach \n in {a,c,h,j,l,o} \node[case,minimum size=.42cm] at (\n) {};
    \foreach \n in {b,e,f,i,m,p} \node[ctrl,minimum size=.42cm] at (\n) {};
    \foreach \n in {d,g,k,n} \node[unk,minimum size=.42cm] at (\n) {};
    \draw[Img,line width=1.1pt] (g) circle (.34);
    \draw[Non,line width=1.1pt] (g) circle (.46);
  \end{scope}
  % 左侧：标题
  \node[anchor=north west,inner sep=0,font=\fontsize{8.6}{10}\selectfont\bfseries,text=NavyMid] at (0,7.3) {《图机器学习》课程汇报};
  \node[anchor=north west,inner sep=0,font=\fontsize{36}{38}\selectfont\bfseries,text=Navy] at (-.05,6.55) {MM-GTUNets};
  \node[anchor=north west,inner sep=0,align=left,font=\fontsize{18}{26}\selectfont\bfseries,text=Navy] at (0,4.95)
    {面向脑疾病预测的\\统一多模态图深度学习};
  \node[anchor=north west,inner sep=0,align=left,text width=7.2cm,font=\fontsize{8.6}{12}\selectfont,text=Muted] at (0,3.05)
    {Unified Multi-Modal Graph Deep Learning\\for Brain Disorders Prediction};
  \node[anchor=north west,inner sep=0,font=\fontsize{8.6}{12}\selectfont,text=Ink] at (0,1.95)
    {Cai et al.\quad ·\quad IEEE Transactions on Medical Imaging, 2025};
  \draw[rule] (0,1.2) -- (6.4,1.2);
  \node[anchor=north west,inner sep=0,font=\fsmall,text=Muted] at (0,.85) {汇报人\quad 〔姓名〕\qquad 〔学号〕};
  \node[anchor=north west,inner sep=0,font=\fsmall,text=Muted] at (0,.3) {西北工业大学};
\end{canvas}
\begin{tikzpicture}[remember picture,overlay]
  \node[anchor=north east,inner sep=0] at ([xshift=-.9cm,yshift=-.55cm]current page.north east)
    {\includegraphics[height=1.05cm]{assets/nwpu-crest.png}};
\end{tikzpicture}
\end{frame}

% ============================================================
%  02 任务
% ============================================================
\setpart{1}\kick{问题}
\begin{frame}{从脑影像到人群图}
\begin{canvas}
  % ---------- 左：单个受试者 ----------
  \node[head] at (0,6.45) {每位受试者};
  \node[anchor=north west,inner sep=0,font=\fsmall\bfseries,text=Img] at (0,5.7) {影像\enspace rs-fMRI};
  \begin{scope}[shift={(0,3.95)}]
    \foreach \b/\p/\c in {1.05/0/100, .62/70/75, .2/150/50}
      \draw[Img!\c,line width=.75pt,smooth] plot[domain=0:1.55,samples=40]
        (\x,{\b+.15*sin(\x*360*1.2+\p)+.08*sin(\x*360*3.3+2*\p)});
  \end{scope}
  \draw[arr] (1.75,4.6) -- (2.2,4.6);
  \begin{scope}[shift={(2.35,3.9)}]
    \foreach \i in {0,...,6} \foreach \j in {0,...,6} {
      \pgfmathsetmacro{\v}{\i==\j ? 95 : 15+60*abs(sin(((\i+\j)*53+abs(\i-\j)*31)))}
      \fill[Img!\v] (\j*.2,1.2-\i*.2) rectangle ++(.2,.2);
    }
  \end{scope}
  \draw[arr] (3.95,4.6) -- (4.4,4.6);
  \begin{scope}[shift={(4.55,4.44)}]
    \foreach \k/\v in {0/70,1/30,2/55,3/85,4/40,5/20,6/65,7/45,8/75}
      \fill[Img!\v] (\k*.19,0) rectangle ++(.17,.32);
  \end{scope}
  \node[note] at (0,3.65) {116 个脑区};
  \node[note] at (2.35,3.65) {功能连接矩阵};
  \node[note] at (4.55,3.65) {6670 维};

  \node[anchor=north west,inner sep=0,font=\fsmall\bfseries,text=Non] at (0,2.65) {非影像\enspace 表型信息};
  \foreach \k/\t in {0/性别,1/年龄,2/采集站点}
    \node[nonbox,anchor=west,minimum height=.5cm,font=\fnote] at ({\k*1.05},1.75) {\t};
  \draw[arr] (3.95,1.75) -- (4.4,1.75);
  \begin{scope}[shift={(4.55,1.59)}]
    \foreach \k/\v in {0/70,1/40,2/90}
      \fill[Non!\v] (\k*.19,0) rectangle ++(.17,.32);
  \end{scope}
  \node[note] at (4.55,1.27) {3 维};
  \draw[Navy,line width=.7pt,dashed,->] (6.4,3.2) to[out=0,in=180] (8.15,3.35);
  \node[anchor=south,inner sep=0,font=\fnote,text=Navy] at (7.1,3.38) {节点特征};

  % ---------- 右：人群图 ----------
  \draw[rule] (6.95,1.15) -- (6.95,6.45);
  \node[head] at (7.45,6.45) {人群图};
  \foreach \n/\x/\y in {n1/8.7/5.0, n2/9.9/5.55, n3/8.5/3.35, n4/9.9/4.2, n5/11.1/5.1, n6/10.6/2.9,
                        n7/12.1/4.0, n8/13.0/5.3, n9/9.3/1.95, n10/11.7/1.85, n11/13.2/2.65, n12/13.75/4.15}
    \coordinate (\n) at (\x,\y);
  \foreach \u/\v/\w/\o in {n1/n2/1.4/.55, n1/n4/.7/.4, n2/n5/.9/.45, n3/n4/1.3/.55, n4/n6/.6/.35,
                           n4/n5/1.0/.5, n5/n7/.8/.4, n6/n9/1.2/.5, n6/n10/.7/.4, n7/n8/1.3/.55,
                           n7/n10/.6/.35, n8/n12/.9/.45, n11/n12/1.1/.5, n10/n11/.9/.45, n3/n9/.8/.4,
                           n7/n11/.5/.3, n2/n4/.6/.35, n1/n3/.5/.3}
    \draw[Navy,opacity=\o,line width=\w pt] (\u) -- (\v);
  \foreach \n in {n1,n2,n6,n8,n11} \node[case] at (\n) {};
  \foreach \n in {n3,n5,n7,n9,n12} \node[ctrl] at (\n) {};
  \foreach \n in {n4,n10} \node[unk] at (\n) {};
  \draw[Img,line width=1pt] (n3) circle (.27);
  \draw[Non,line width=1pt] (n3) circle (.37);
  \node[case,minimum size=.24cm] at (10.35,6.33) {}; \node[anchor=west,inner sep=0,font=\fnote] at (10.55,6.33) {患者};
  \node[ctrl,minimum size=.24cm] at (11.5,6.33) {}; \node[anchor=west,inner sep=0,font=\fnote] at (11.7,6.33) {健康对照};
  \node[unk,minimum size=.24cm] at (13.1,6.33) {}; \node[anchor=west,inner sep=0,font=\fnote] at (13.3,6.33) {待预测};

  \takeaway{一个节点是一位受试者，诊断转化为人群图上的节点分类}
\end{canvas}
\end{frame}

% ============================================================
%  03 挑战与对策
% ============================================================
\begin{frame}{现有方法的三个不足}
\begin{canvas}
  \foreach \k/\nn/\prob/\desc/\mtag/\sol in {
    0/01/非影像数据利用不足/{表型信息只用来计算边权}/MRRL/{VAE 重建表型表示\\奖励机制学习属性权重},
    1/02/关键节点被忽视/{全图节点被同等对待}/GTUNet/{图池化筛选关键节点\\图 Transformer 获取全局信息},
    2/03/跨模态交互不足/{两类数据只在节点与边上交互}/ACMGL/{双通道分别编码\\注意力融合两类表示}} {
    \begin{scope}[shift={(\k*4.9,0)}]
      \node[anchor=north west,inner sep=0,font=\fontsize{20}{22}\selectfont\bfseries,text=Faint] at (0,6.45) {\nn};
      \node[head,font=\fontsize{11.5}{15}\selectfont\bfseries] at (0,5.5) {\prob};
      \node[txt,font=\fsmall,text=Muted] at (0,4.85) {\desc};
      \fill[NavyPale,rounded corners=3pt] (0,1.8) rectangle (4.4,3.95);
      \fill[Navy] (0,2.25) rectangle (.06,3.5);
      \node[anchor=north west,inner sep=0,font=\fnote\bfseries,text=NavyMid] at (.32,3.62) {\mtag};
      \node[anchor=north west,inner sep=0,align=left,font=\fsmall\bfseries,text=Navy] at (.32,3.12) {\sol};
    \end{scope}
  }
  \takeaway{MM-GTUNets：构图、编码与融合在一个框架内端到端训练}
\end{canvas}
\end{frame}

% ============================================================
%  04 整体框架
% ============================================================
\setpart{2}\kick{方法}\setstage{0}
\begin{frame}{整体框架}
\begin{canvas}
  \fill[Pale,rounded corners=4pt] (2.1,.75) rectangle (6.85,6.45);
  \fill[Pale,rounded corners=4pt] (7.1,.75) rectangle (11.75,6.45);
  \fill[Pale,rounded corners=4pt] (12.0,.75) rectangle (14.2,6.45);
  \node[anchor=north west,inner sep=0,font=\fnote\bfseries,text=Navy] at (2.3,6.25) {模态奖励表示学习\enspace\textcolor{Muted}{MRRL}};
  \node[anchor=north west,inner sep=0,font=\fnote\bfseries,text=Navy] at (7.3,6.25) {自适应跨模态图学习\enspace\textcolor{Muted}{ACMGL}};
  \node[anchor=north west,inner sep=0,font=\fnote\bfseries,text=Navy] at (12.2,6.25) {分类};

  % 三个核心环节：与方法各页右上角的导航一致
  \node[pillon] at (5.75,5.62) {构图};
  \node[pillon] at (8.55,5.62) {编码};
  \node[pillon] at (10.75,5.62) {融合};

  \node[imgbox,text width=1.45cm] (in1) at (.9,4.85) {rs-fMRI\\功能连接};
  \node[box,text width=1.45cm] (inl) at (.9,3.3) {已知标签};
  \node[nonbox,text width=1.45cm] (in2) at (.9,1.75) {性别 · 年龄\\采集站点};

  \node[imgbox,text width=1.6cm] (rfe) at (3.25,4.85) {RFE 降维\\500 维};
  \node[nonbox,text width=1.6cm] (vae) at (3.25,1.75) {VAE 重建\\500 维};
  \node[box,draw=Navy,line width=1pt,text width=1.6cm] (amrs) at (3.25,3.3) {AMRS\\奖励构图};
  \node[navbox,draw=Navy,line width=1pt,text width=1.25cm,minimum height=1.25cm] (A) at (5.75,3.3) {};
  \begin{scope}[shift={(5.75,3.38)},scale=.33]
    \foreach \a/\b/\c/\d in {-1/.8/.2/1.1, .2/1.1/1/.2, -1/.8/-.6/-.4, -.6/-.4/1/.2, -.6/-.4/.5/-.9, .2/1.1/-.6/-.4}
      \draw[Navy,line width=.6pt] (\a,\b) -- (\c,\d);
    \foreach \a/\b in {-1/.8, .2/1.1, 1/.2, -.6/-.4, .5/-.9} \fill[Navy] (\a,\b) circle (.2);
  \end{scope}
  \node[anchor=north,inner sep=0,font=\fnote\bfseries,text=Navy] at (5.75,2.52) {人群图 $A$};
  \draw[imgarr] (in1) -- (rfe);
  \draw[nonarr] (in2) -- (vae);
  \draw[->,line width=1pt,draw=Muted] (inl) -- (amrs);
  \draw[nonarr] (in2.north east) to[out=50,in=-140] ([xshift=.3cm]amrs.south west);
  \draw[imgarr,line width=.7pt] (rfe.south east) to[out=-30,in=150] (A.north west);
  \draw[nonarr,line width=.7pt] (vae.north east) to[out=30,in=-150] (A.south west);
  \draw[navarr,line width=.8pt] (amrs) -- node[above,font=\ftiny,text=Navy] {$C$} (A);

  \node[imgbox,draw=Img,line width=1pt,text width=1.75cm] (g1) at (8.55,4.85) {GTUNet\\影像通道};
  \node[nonbox,draw=Non,line width=1pt,text width=1.75cm] (g2) at (8.55,1.75) {GTUNet\\非影像通道};
  \node[navbox,draw=Navy,line width=1pt,text width=1.3cm] (fu) at (10.75,3.3) {注意力\\融合};
  \draw[imgarr] (rfe) -- (g1);
  \draw[nonarr] (vae) -- (g2);
  \draw[Navy,line width=.8pt,dashed,->] (A.east) -| (g1.south);
  \draw[Navy,line width=.8pt,dashed,->] (A.east) -| (g2.north);
  \draw[imgarr] (g1.east) to[out=0,in=90] (fu.north);
  \draw[nonarr] (g2.east) to[out=0,in=-90] (fu.south);

  \node[box,text width=1.2cm] (mlp) at (13.1,3.3) {MLP};
  \node[anchor=north,inner sep=0,align=center,font=\fnote,text=Ink] (y) at (13.1,2.35) {患者 / 健康对照};
  \draw[navarr] (fu) -- node[above,font=\ftiny,text=Navy] {$Z$} (mlp);
  \draw[arr] (mlp) -- (y);
\end{canvas}
\end{frame}

% ============================================================
%  05 AMRS
% ============================================================
\setstage{1}
\begin{frame}{奖励机制学习属性权重}
\begin{canvas}
  \node[head] at (0,6.45) {以“采集站点”为例};
  \node[case,minimum size=.62cm,font=\fsmall\bfseries,text=white] (s1) at (1.0,3.6) {1};
  \node[case,minimum size=.62cm,font=\fsmall\bfseries,text=white] (s2) at (3.55,5.1) {2};
  \node[ctrl,minimum size=.62cm,font=\fsmall\bfseries,text=Case] (s3) at (3.55,3.6) {3};
  \node[unk,minimum size=.62cm,font=\fsmall\bfseries,text=Muted] (s4) at (3.55,2.1) {4};
  \draw[Rew,line width=1.6pt] (s1) -- (s2);
  \draw[Pen,line width=1.6pt] (s1) -- (s3);
  \draw[Faint,line width=1.3pt,dashed] (s1) -- (s4);
  \node[anchor=west,inner sep=0,align=left,font=\fsmall] at (4.15,5.1) {\textcolor{Rew}{\bfseries 奖励}\\[-1pt]\textcolor{Muted}{\fnote 同站点，同类别}};
  \node[anchor=west,inner sep=0,align=left,font=\fsmall] at (4.15,3.6) {\textcolor{Pen}{\bfseries 惩罚}\\[-1pt]\textcolor{Muted}{\fnote 同站点，异类别}};
  \node[anchor=west,inner sep=0,align=left,font=\fsmall] at (4.15,2.1) {\textcolor{Muted}{\bfseries 激励}\\[-1pt]\textcolor{Muted}{\fnote 同站点，标签未知}};
  \node[case,minimum size=.24cm] at (.12,1.15) {}; \node[anchor=west,inner sep=0,font=\fnote] at (.32,1.15) {患者};
  \node[ctrl,minimum size=.24cm] at (1.3,1.15) {}; \node[anchor=west,inner sep=0,font=\fnote] at (1.5,1.15) {健康对照};
  \node[unk,minimum size=.24cm] at (2.95,1.15) {}; \node[anchor=west,inner sep=0,font=\fnote] at (3.15,1.15) {待预测};

  \draw[rule] (6.85,1.0) -- (6.85,6.45);
  \node[head] at (7.35,6.45) {非影像亲和度};
  \node[anchor=north west,inner sep=0,font=\fontsize{11}{14}\selectfont] at (7.35,5.6)
    {$C_{ij}=\sigma\Big(\sum_{u=1}^{v}\alpha_u\big(\beta_r^{u}R_{ij}+\beta_p^{u}P_{ij}+\beta_m^{u}M_{ij}\big)\Big)$};
  \foreach \y/\t/\d in {4.05/{$R,\,P,\,M$}/{奖励、惩罚、激励计数},
                        3.25/{$\beta$}/{奖励与激励为正，惩罚为负},
                        2.45/{$\alpha_u$}/{属性权重，训练中学习}} {
    \node[anchor=west,inner sep=0,font=\fbody,text=Navy] at (7.35,\y) {\t};
    \node[anchor=west,inner sep=0,font=\fsmall] at (8.75,\y) {\d};
  }
  \node[anchor=west,inner sep=0,font=\fsmall] at (7.35,1.45) {借鉴 \textcolor{Navy}{\bfseries Q-learning}：以奖励价值 $Q$ 引导 $\alpha$ 的学习};

  \takeaway{同属性、同类别的受试者连接更强}
\end{canvas}
\end{frame}

% ============================================================
%  06 人群图构建
% ============================================================
\begin{frame}{自适应人群图}
\begin{canvas}
  \node[head] at (0,6.45) {模态对齐};
  \node[imgbox,text width=1.25cm] (a1) at (.75,5.0) {影像\\6670 维};
  \node[imgbox,text width=1.0cm] (a2) at (3.35,5.0) {500 维};
  \node[nonbox,text width=1.25cm] (b1) at (.75,3.0) {非影像\\3 维};
  \node[nonbox,text width=1.0cm] (b2) at (3.35,3.0) {500 维};
  \draw[imgarr] (a1) -- node[above,font=\fnote\bfseries,text=Img] {RFE} (a2);
  \draw[nonarr] (b1) -- node[above,font=\fnote\bfseries,text=Non] {VAE} (b2);
  \node[navbox,text width=1.15cm] (cat) at (5.6,4.0) {拼接\\$\hat{x}_i$};
  \draw[imgarr] (a2.east) to[out=0,in=150] (cat.north west);
  \draw[nonarr] (b2.east) to[out=0,in=-150] (cat.south west);
  \node[anchor=west,inner sep=0,font=\fsmall] at (0,1.45) {两类特征对齐到\textcolor{Navy}{\bfseries 同一维度}后拼接};

  \draw[rule] (6.85,1.0) -- (6.85,6.45);
  \node[head] at (7.35,6.45) {边权};
  \node[anchor=north west,inner sep=0,font=\fontsize{12.5}{15}\selectfont] at (7.35,5.7)
    {$A_{ij}=\mathrm{Sim}(\hat{x}_i,\hat{x}_j)\odot C_{ij}$};
  \foreach \y/\w/\o/\c/\a/\st in {3.85/3.2/.85/0.90/0.72/case, 2.45/1.0/.45/0.30/0.24/ctrl} {
    \draw[Navy,opacity=\o,line width=\w pt] (7.65,\y) -- (9.05,\y);
    \node[\st,minimum size=.4cm] at (7.65,\y) {};
    \node[case,minimum size=.4cm] at (9.05,\y) {};
    \node[anchor=west,inner sep=0,font=\fsmall] at (9.6,\y)
      {0.80\enspace$\times$\enspace\c\enspace$=$\enspace\textcolor{Navy}{\bfseries\a}};
  }
  \node[anchor=south west,inner sep=0,font=\fnote,text=Muted] at (9.6,4.45) {Sim\enspace$\times$\enspace C};
  \node[anchor=west,inner sep=0,font=\fsmall] at (7.35,1.45) {特征同样相似时，\textcolor{Navy}{\bfseries 属性亲和度}决定边的强弱};

  \takeaway{边权 = 特征相似度 × 属性亲和度}
\end{canvas}
\end{frame}

% ============================================================
%  07 图 Transformer 层
% ============================================================
\setstage{2}
\begin{frame}{带边特征的图 Transformer}
\begin{canvas}
  \node[circle,fill=Navy,minimum size=.72cm,inner sep=0,font=\fsmall\bfseries,text=white] (i) at (3.9,3.8) {$i$};
  \foreach \n/\x/\y/\w in {j1/1.1/5.6/2.6, j2/0.45/3.45/1.6, j3/2.05/1.85/0.9} {
    \node[circle,draw=Navy,line width=.8pt,fill=white,minimum size=.56cm,inner sep=0] (\n) at (\x,\y) {};
    \draw[->,Navy,line width=\w pt,shorten >=1pt] (\n) -- (i);
  }
  \node[font=\fnote\bfseries,text=Navy,fill=white,inner sep=1.5pt] at ($(j1)!.5!(i)+(0.2,.2)$) {0.52};
  \node[font=\fnote\bfseries,text=Navy,fill=white,inner sep=1.5pt] at ($(j2)!.5!(i)+(0,.22)$) {0.31};
  \node[font=\fnote\bfseries,text=Navy,fill=white,inner sep=1.5pt] at ($(j3)!.5!(i)+(.24,-.12)$) {0.17};
  \node[anchor=west,font=\fnote,text=Muted,inner sep=2pt] at (j1.east) {$k_j,\ v_j$};
  \node[anchor=west,font=\fnote,text=Muted,inner sep=3pt] at (i.east) {$q_i$};
  \node[anchor=west,inner sep=0,font=\fsmall] at (0,1.15) {箭头粗细表示注意力权重 $\alpha_{ij}$};

  \draw[rule] (6.2,1.0) -- (6.2,6.45);
  \node[anchor=north west,inner sep=0,font=\fsmall\bfseries,text=NavyMid] at (6.7,6.45) {注意力};
  \node[anchor=north west,inner sep=0,font=\fbody] at (6.7,6.0)
    {$\alpha_{ij}=\dfrac{\exp\big(q_i^{\mathrm{T}}(k_j+e_{ij})/\sqrt{d}\big)}{\sum_{u\in\mathcal{N}(i)}\exp\big(q_i^{\mathrm{T}}(k_u+e_{iu})/\sqrt{d}\big)}$};
  \node[anchor=north west,inner sep=0,font=\fsmall\bfseries,text=NavyMid] at (6.7,4.5) {消息};
  \node[anchor=north west,inner sep=0,font=\fbody] at (6.7,4.05)
    {$\bar{h}_i=\sum_{j\in\mathcal{N}(i)}\alpha_{ij}\,(v_j+e_{ij})$};
  \node[anchor=north west,inner sep=0,font=\fsmall\bfseries,text=NavyMid] at (6.7,3.05) {门控残差};
  \node[anchor=north west,inner sep=0,font=\fbody,align=left] at (6.7,2.6)
    {$\gamma_i=\sigma\big(W_g[\,\bar{h}_i;\,r_i;\,\bar{h}_i-r_i\,]\big)$\\[3pt]
     $h_i'=\mathrm{ReLU}\big(\mathrm{LN}\big((1-\gamma_i)\,\bar{h}_i+\gamma_i\,r_i\big)\big)$};

  \takeaway{边特征 $e_{ij}$ 同时参与注意力打分和消息传递}
\end{canvas}
\end{frame}

% ============================================================
%  08 GTUNet
% ============================================================
\begin{frame}{GTUNet：U 形图 Transformer}
\begin{canvas}
  \foreach \n/\x/\y/\c in {b1/0.85/5.6/N, b2/3.35/4.6/0.8N, b3/5.85/3.6/0.64N, b4/8.35/4.6/0.8N, b5/10.85/5.6/N} {
    \node[navbox,minimum width=1.5cm,minimum height=.74cm,align=center,inner ysep=2pt] (\n) at (\x,\y)
      {{\fsmall\bfseries\color{Navy}GT 层}\\[-2pt]{\ftiny\color{Muted}\c}};
  }
  \draw[navarr] (b1.south) |- (b2.west);
  \draw[navarr] (b2.south) |- (b3.west);
  \draw[navarr] (b3.east) -| (b4.south);
  \draw[navarr] (b4.east) -| (b5.south);
  \node[anchor=west,inner sep=0,font=\fnote\bfseries,text=Navy] at (.97,4.85) {gPool};
  \node[anchor=west,inner sep=0,font=\fnote\bfseries,text=Navy] at (3.47,3.85) {gPool};
  \node[anchor=east,inner sep=0,font=\fnote\bfseries,text=Navy] at (8.23,3.85) {gUnpool};
  \node[anchor=east,inner sep=0,font=\fnote\bfseries,text=Navy] at (10.73,4.85) {gUnpool};
  \draw[Faint,line width=.8pt,dashed,->] (b1.north) -- ++(0,.25) -| (b5.north);
  \draw[Faint,line width=.8pt,dashed,->] (b2.north) -- ++(0,.3) -| (b4.north);
  \node[anchor=north west,inner sep=0,align=left,font=\fsmall] at (12.0,6.1)
    {\textcolor{Muted}{编码深度}\\影像 2\\非影像 3\\[5pt]\textcolor{Muted}{池化比例}\\0.8};

  \draw[rule] (0,2.75) -- (14.2,2.75);
  \node[anchor=north west,inner sep=0,font=\fsmall\bfseries,text=Navy] at (0,2.45) {gPool\enspace 按得分保留前 80\% 的节点};
  \foreach \k/\s/\keep in {0/0.9/1, 1/0.2/0, 2/0.7/1, 3/0.5/1, 4/0.8/1} {
    \ifnum\keep=1
      \node[circle,fill=Navy,minimum size=.42cm,inner sep=0] at (\k*.85+.25,1.55) {};
      \node[font=\fnote,text=Navy,inner sep=0] at (\k*.85+.25,1.0) {\s};
    \else
      \node[circle,draw=Muted,dashed,minimum size=.42cm,inner sep=0] at (\k*.85+.25,1.55) {};
      \node[font=\fnote,text=Muted,inner sep=0] at (\k*.85+.25,1.0) {\s};
    \fi
  }
  \node[anchor=north west,inner sep=0,font=\fsmall\bfseries,text=Navy] at (7.35,2.45) {gUnpool\enspace 按索引写回池化前的特征};
  \foreach \k/\keep in {0/1, 1/0, 2/1, 3/1, 4/1} {
    \ifnum\keep=1
      \node[circle,fill=Navy,minimum size=.42cm,inner sep=0] at (7.6+\k*.85,1.55) {};
    \else
      \node[circle,fill=NavyPale,draw=Navy!45,line width=.8pt,minimum size=.42cm,inner sep=0] at (7.6+\k*.85,1.55) {};
    \fi
  }
  \fill[Navy] (11.85,1.78) circle (.07); \node[anchor=west,inner sep=0,font=\fnote,text=Muted] at (12.05,1.78) {池化后更新};
  \fill[NavyPale] (11.85,1.32) circle (.07); \draw[Navy!45,line width=.6pt] (11.85,1.32) circle (.07);
  \node[anchor=west,inner sep=0,font=\fnote,text=Muted] at (12.05,1.32) {池化前特征};
  \takeaway{池化聚焦关键节点，反池化恢复完整人群图}
\end{canvas}
\end{frame}

% ============================================================
%  09 融合与训练目标
% ============================================================
\setstage{3}
\begin{frame}{多模态注意力融合}
\begin{canvas}
  \node[imgbox,minimum width=1.25cm] (zi) at (.7,5.85) {$Z^{img}_{sp}$};
  \node[nonbox,minimum width=1.25cm] (zn) at (.7,3.95) {$Z^{non}_{sp}$};
  \node[navbox,minimum width=1.25cm] (zs) at (2.6,4.9) {$Z_{sh}$};
  \draw[imgarr,line width=.8pt] (zi.east) to[out=0,in=150] (zs.north west);
  \draw[nonarr,line width=.8pt] (zn.east) to[out=0,in=-150] (zs.south west);
  \foreach \r/\y/\c/\t in {zi/5.85/Img/{\tau^{img}_{sp}}, zs/4.9/Navy/{\tau_{sh}}, zn/3.95/Non/{\tau^{non}_{sp}}} {
    \node[circle,draw=\c,line width=.9pt,fill=white,minimum size=.46cm,inner sep=0,font=\fsmall,text=\c] (g\r) at (4.75,\y) {$\odot$};
    \node[anchor=south east,inner sep=1.5pt,font=\fnote,text=\c] at (4.45,\y) {$\t$};
    \draw[->,line width=.9pt,draw=\c] (\r.east) -- (g\r.west);
  }
  \node[circle,draw=Navy,line width=.9pt,minimum size=.5cm,inner sep=0,font=\fsmall,text=Navy] (sum) at (6.2,4.9) {$+$};
  \foreach \r in {zi,zs,zn} \draw[->,line width=.9pt,draw=Faint] (g\r.east) to[out=0,in=180] (sum.west);
  \node[navbox,minimum width=1.0cm,font=\fsmall\bfseries,text=Navy] (z) at (7.4,4.9) {$Z$};
  \draw[navarr] (sum) -- (z);

  \node[anchor=north west,inner sep=0,align=left,font=\fsmall] at (8.55,6.15)
    {$Z_{sh}=\tfrac12\big(Z^{img}_{sp}+Z^{non}_{sp}\big)$\\[7pt]
     $\tau_{\ast}=\tanh\big(W_{\ast}Z_{\ast}+B_{\ast}\big)$\\[7pt]
     $Z=\tau_{sh}\odot Z_{sh}+\tau^{img}_{sp}\odot Z^{img}_{sp}+\tau^{non}_{sp}\odot Z^{non}_{sp}$};

  \draw[rule] (0,3.05) -- (14.2,3.05);
  \node[anchor=north west,inner sep=0,font=\fsmall\bfseries,text=Navy] at (0,2.65) {训练目标};
  \node[anchor=north west,inner sep=0,font=\fontsize{11.5}{14}\selectfont] at (2.0,2.7)
    {$\mathcal{L}_{total}=\mathcal{L}_{ce}+\omega_{img}\,\mathcal{L}^{img}_{g}+\omega_{non}\big(\mathcal{L}^{non}_{g}+\eta\,\mathcal{L}_{r}\big)$};
  \foreach \x/\a/\b in {2.0/{$\mathcal{L}_{ce}$}/{分类损失},
                        4.75/{$\mathcal{L}_{g}$}/{图正则：平滑与度},
                        8.15/{$\mathcal{L}_{r}$}/{奖励正则},
                        10.6/{$\omega$}/{由 $\tau$ 算出的模态权重}} {
    \node[anchor=west,inner sep=0,font=\fsmall] at (\x,1.5) {\a\enspace\textcolor{Muted}{\b}};
  }
  \takeaway{联合表示同时保留模态特有信息与共享信息}
\end{canvas}
\end{frame}

% ============================================================
%  10 实验设置
% ============================================================
\setpart{3}\kick{实验}\setstage{0}
\begin{frame}{实验设置}
\begin{canvas}
  \node[head] at (0,6.45) {数据集};
  \foreach \y/\name/\desc/\pn/\pl/\hn in {
      5.7/ABIDE/{自闭症谱系障碍\enspace·\enspace 871 人\enspace·\enspace 20 个站点}/403/ASD/468,
      3.45/ADHD-200/{注意缺陷多动障碍\enspace·\enspace 582 人\enspace·\enspace 4 个站点}/218/ADHD/364} {
    \node[anchor=north west,inner sep=0,font=\fontsize{11.5}{13}\selectfont\bfseries,text=Navy] at (0,\y) {\name};
    \node[anchor=north west,inner sep=0,font=\fnote,text=Muted] at (0,\y-.5) {\desc};
    \pgfmathsetmacro{\wp}{6.5*\pn/(\pn+\hn)}
    \fill[Case] (0,\y-1.5) rectangle (\wp,\y-1.05);
    \fill[Ctrl] (\wp,\y-1.5) rectangle (6.5,\y-1.05);
    \node[anchor=west,inner sep=0,font=\fnote\bfseries,text=white] at (.15,\y-1.275) {\pl\enspace\pn};
    \node[anchor=east,inner sep=0,font=\fnote\bfseries,text=Case] at (6.35,\y-1.275) {健康对照\enspace\hn};
  }

  \draw[rule] (7.0,1.0) -- (7.0,6.45);
  \node[head] at (7.5,6.45) {评估方式};
  \node[anchor=north west,inner sep=0,align=left,font=\fsmall] at (7.5,5.7)
    {\nav{十折分层交叉验证}\\[2pt]训练、验证、测试按 8\,:\,1\,:\,1 划分};
  \node[anchor=north west,inner sep=0,align=left,font=\fsmall] at (7.5,4.55)
    {\nav{评价指标}\\[2pt]ACC\enspace·\enspace SEN\enspace·\enspace SPE\enspace·\enspace AUC};
  \node[anchor=north west,inner sep=0,font=\fsmall] at (7.5,3.4) {\nav{对比方法}};
  \foreach \y/\g/\m in {2.75/传统/{SVM、MLP}, 2.2/脑图/{Brain-GNN、DGCN、AL-NEGAT、A-GCL}, 1.65/人群图/{Pop-GCN、GATE、EV-GCN、MMGL}} {
    \node[anchor=west,inner sep=0,font=\fnote,text=Muted] at (7.5,\y) {\g};
    \node[anchor=west,inner sep=0,font=\fsmall] at (8.55,\y) {\m};
  }

  \takeaway{两项二分类任务：ASD 与健康对照，ADHD 与健康对照}
\end{canvas}
\end{frame}

% ============================================================
%  11 主要结果
% ============================================================
\begin{frame}{分类性能对比}
\begin{canvas}
  \def\rows{传统/SVM/5.6/66.02/66.48, /MLP/5.22/72.69/75.28,
            脑图/Brain-GNN/4.7/69.76/65.26, /DGCN/4.32/71.83/68.81, /AL-NEGAT/3.94/73.17/69.15, /A-GCL/3.56/79.04/80.11,
            人群图/Pop-GCN/3.04/68.43/75.45, /GATE/2.66/74.65/72.20, /EV-GCN/2.28/80.95/80.95, /MMGL/1.9/80.34/77.59}
  \fill[NavyPale,rounded corners=2pt] (0,1.12) rectangle (11.3,1.58);
  \foreach \px/\name in {2.75/ABIDE, 7.25/ADHD-200} {
    \node[anchor=south west,inner sep=0,font=\fsmall\bfseries,text=Navy] at (\px,6.25) {\name\enspace\textcolor{Muted}{\fnote ACC (\%)}};
    \foreach \t in {60,70,80} {
      \pgfmathsetmacro{\xx}{\px+(\t-60)/25*3.9}
      \draw[Line,line width=.5pt] (\xx,1.15) -- (\xx,5.85);
      \node[anchor=south,inner sep=0,font=\ftiny,text=Muted] at (\xx,5.9) {\t};
    }
  }
  \foreach \g/\m/\y/\a/\b in \rows {
    \node[anchor=west,inner sep=0,font=\fnote,text=Muted] at (0,\y) {\g};
    \node[anchor=east,inner sep=0,font=\fnote,text=Ink] at (2.55,\y) {\m};
    \pgfmathsetmacro{\xa}{2.75+(\a-60)/25*3.9}
    \pgfmathsetmacro{\xb}{7.25+(\b-60)/25*3.9}
    \fill[Faint] (\xa,\y) circle (.075);
    \fill[Faint] (\xb,\y) circle (.075);
  }
  \node[anchor=west,inner sep=0,font=\ftiny,text=Muted] at ({2.75+(80.95-60)/25*3.9+.14},2.28) {80.95};
  \node[anchor=west,inner sep=0,font=\ftiny,text=Muted] at ({7.25+(80.95-60)/25*3.9+.14},2.28) {80.95};
  \node[anchor=west,inner sep=0,font=\fnote\bfseries,text=Navy] at (0,1.35) {本文};
  \node[anchor=east,inner sep=0,font=\fnote\bfseries,text=Navy] at (2.55,1.35) {MM-GTUNets};
  \fill[Navy] ({2.75+(82.92-60)/25*3.9},1.35) circle (.11);
  \fill[Navy] ({7.25+(82.68-60)/25*3.9},1.35) circle (.11);
  \node[anchor=west,inner sep=0,font=\fnote\bfseries,text=Navy] at ({2.75+(82.92-60)/25*3.9+.17},1.35) {82.92};
  \node[anchor=west,inner sep=0,font=\fnote\bfseries,text=Navy] at ({7.25+(82.68-60)/25*3.9+.17},1.35) {82.68};
  \foreach \y in {4.96,3.3,1.64} \draw[Line,line width=.4pt] (0,\y) -- (2.55,\y);

  \draw[rule] (11.7,1.15) -- (11.7,6.4);
  \foreach \y/\name/\acc/\da/\auc/\db in {6.4/ABIDE/82.92/1.97/88.21/5.35, 3.65/ADHD-200/82.68/1.73/90.71/3.56} {
    \node[anchor=north west,inner sep=0,font=\fnote\bfseries,text=Muted] at (12.05,\y) {\name};
    \node[anchor=north west,inner sep=0,font=\fbig,text=Navy] at (12.0,\y-.42) {\acc{\fontsize{12}{14}\selectfont\%}};
    \node[anchor=north west,inner sep=0,align=left,font=\fnote] at (12.05,\y-1.35)
      {ACC\enspace\textcolor{Navy}{\bfseries +\da}\\AUC\enspace\textcolor{Navy}{\bfseries +\db}};
  }

  \takeaway{ABIDE 四项指标均最优，ADHD-200 的 ACC 与 AUC 最优}
\end{canvas}
\end{frame}

% ============================================================
%  12 消融实验
% ============================================================
\begin{frame}{消融实验}
\begin{canvas}
  \def\rows{构图/Pop-Graph/5.65/77.40/78.03/0, /MCA-Graph/5.29/78.90/81.46/0, /EV-Graph/4.93/80.38/80.80/0, /MRRL/4.57/82.92/82.68/1,
            重建/不重建/4.0/67.06/76.47/0, /MLP/3.64/80.86/79.04/0, /AE/3.28/76.02/77.86/0, /VAE/2.92/82.92/82.68/1,
            编码器/堆叠/2.35/77.06/76.14/0, /残差/1.99/78.78/76.33/0, /级联/1.63/79.70/80.97/0, /GTUNet/1.27/82.92/82.68/1}
  \foreach \y in {4.57,2.92,1.27} \fill[NavyPale,rounded corners=2pt] (0,\y-.17) rectangle (11.3,\y+.17);
  \foreach \px/\name in {2.75/ABIDE, 7.25/ADHD-200} {
    \node[anchor=south west,inner sep=0,font=\fsmall\bfseries,text=Navy] at (\px,6.25) {\name\enspace\textcolor{Muted}{\fnote ACC (\%)}};
    \foreach \t in {70,75,80,85} {
      \pgfmathsetmacro{\xx}{\px+(\t-65)/20*3.9}
      \draw[Line,line width=.5pt] (\xx,1.1) -- (\xx,5.85);
      \node[anchor=south,inner sep=0,font=\ftiny,text=Muted] at (\xx,5.9) {\t};
    }
  }
  \foreach \g/\m/\y/\a/\b/\o in \rows {
    \node[anchor=west,inner sep=0,font=\fnote,text=Muted] at (0,\y) {\g};
    \pgfmathsetmacro{\xa}{2.75+(\a-65)/20*3.9}
    \pgfmathsetmacro{\xb}{7.25+(\b-65)/20*3.9}
    \ifnum\o=1
      \node[anchor=east,inner sep=0,font=\fnote\bfseries,text=Navy] at (2.55,\y) {\m};
      \fill[Navy] (\xa,\y) circle (.1);
      \fill[Navy] (\xb,\y) circle (.1);
      \node[anchor=west,inner sep=0,font=\ftiny\bfseries,text=Navy] at (\xa+.15,\y) {\a};
      \node[anchor=west,inner sep=0,font=\ftiny\bfseries,text=Navy] at (\xb+.15,\y) {\b};
    \else
      \node[anchor=east,inner sep=0,font=\fnote,text=Ink] at (2.55,\y) {\m};
      \fill[Faint] (\xa,\y) circle (.075);
      \fill[Faint] (\xb,\y) circle (.075);
    \fi
  }
  \foreach \y in {4.285,2.635} \draw[Line,line width=.4pt] (0,\y) -- (2.55,\y);

  \draw[rule] (11.7,1.1) -- (11.7,6.4);
  \node[anchor=south west,inner sep=0,font=\fsmall\bfseries,text=Navy] at (11.95,6.25) {较次优提升};
  \node[anchor=south west,inner sep=0,font=\ftiny,text=Muted] at (11.95,5.9) {ABIDE};
  \node[anchor=south west,inner sep=0,font=\ftiny,text=Muted] at (13.0,5.9) {ADHD-200};
  \foreach \y/\a/\b in {5.11/2.54/1.22, 3.46/2.06/3.64, 1.81/3.22/1.71} {
    \node[anchor=west,inner sep=0,font=\fbody\bfseries,text=Navy] at (11.95,\y) {+\a};
    \node[anchor=west,inner sep=0,font=\fbody\bfseries,text=Navy] at (13.0,\y) {+\b};
  }

  \takeaway{三处设计在两个数据集上均优于次优的替代方案}
\end{canvas}
\end{frame}

% ============================================================
%  13 模态作用
% ============================================================
\begin{frame}{影像与非影像的作用}
\begin{canvas}
  \node[head] at (0,6.45) {输入消融\enspace\textcolor{Muted}{\fnote ACC (\%)}};
  \foreach \y/\lab/\c/\a/\b in {5.3/仅非影像/Non/52.00/58.07, 4.05/仅影像/Img/76.49/79.58, 2.8/影像 + 非影像/Navy/82.92/82.68} {
    \node[anchor=west,inner sep=0,font=\fsmall\bfseries,text=\c] at (0,\y) {\lab};
    \fill[\c] (2.25,\y+.05) rectangle ({2.25+\a/100*3.9},\y+.33);
    \fill[\c!40] (2.25,\y-.33) rectangle ({2.25+\b/100*3.9},\y-.05);
    \node[anchor=west,inner sep=0,font=\fnote,text=Ink] at ({2.25+\a/100*3.9+.1},\y+.19) {\a};
    \node[anchor=west,inner sep=0,font=\fnote,text=Muted] at ({2.25+\b/100*3.9+.1},\y-.19) {\b};
  }
  \draw[Line,line width=.6pt] (2.25,2.25) -- (2.25,5.85);
  \fill[Ink!70] (2.25,1.62) rectangle ++(.3,.17);
  \node[anchor=west,inner sep=0,font=\fnote] at (2.65,1.705) {ABIDE};
  \fill[Ink!25] (3.75,1.62) rectangle ++(.3,.17);
  \node[anchor=west,inner sep=0,font=\fnote] at (4.15,1.705) {ADHD-200};

  \draw[rule] (7.0,1.0) -- (7.0,6.45);
  \node[head] at (7.5,6.45) {模态贡献};
  \node[imgbox,anchor=west,font=\fnote\bfseries,text=Img,minimum height=.48cm] at (7.5,5.55) {rs-fMRI};
  \node[anchor=north west,inner sep=0,font=\fsmall] at (7.5,5.05) {贡献权重最高，在两个数据集上都占主导};
  \foreach \k/\t in {0/性别,1/年龄,2/采集站点}
    \node[nonbox,anchor=west,font=\fnote\bfseries,text=Non,minimum height=.48cm] at ({7.5+\k*1.05},3.95) {\t};
  \node[anchor=north west,inner sep=0,font=\fsmall] at (7.5,3.45) {三类表型贡献相近：性别影响更大，年龄更稳定};
  \node[nonbox,anchor=west,font=\fnote\bfseries,text=Non,minimum height=.48cm] at (7.5,2.35) {年龄};
  \node[anchor=north west,inner sep=0,font=\fsmall] at (7.5,1.85) {只加入一种表型时，加入年龄的效果最好};

  \takeaway{影像起主导作用，非影像信息带来稳定补充}
\end{canvas}
\end{frame}

% ============================================================
%  14 总结
% ============================================================
\setpart{4}\kick{总结}
\begin{frame}{总结}
\begin{canvas}
  \node[head] at (0,6.45) {核心贡献};
  \foreach \y/\tg/\md/\tx in {5.45/构图/MRRL/{奖励机制学习属性权重，自适应构建人群图},
                              4.15/编码/GTUNet/{图 Transformer 结合图池化，兼顾全局与关键节点},
                              2.85/融合/ACMGL/{注意力融合模态特有与共享表示}} {
    \node[pillon,anchor=west,minimum width=.95cm] at (0,\y) {\tg};
    \node[anchor=north west,inner sep=0,align=left,font=\fsmall] at (1.3,\y+.3)
      {\textcolor{NavyMid}{\bfseries\md}\\[1pt]\tx};
  }
  \node[anchor=west,inner sep=0,font=\fsmall] at (0,1.5)
    {\textcolor{Muted}{准确率}\quad\nav{ABIDE 82.92\%}\qquad\nav{ADHD-200 82.68\%}};

  \draw[rule] (8.6,1.0) -- (8.6,6.45);
  \node[head] at (9.1,6.45) {局限};
  \foreach \y/\t in {5.55/{转导式学习，新受试者需重新建图},
                     4.65/{仅验证二分类与两类模态},
                     3.75/{显存随图规模增长},
                     2.85/{未利用低质量或缺失模态数据}} {
    \fill[Muted] (9.18,\y-.2) circle (.05);
    \node[anchor=north west,inner sep=0,text width=4.8cm,font=\fsmall] at (9.4,\y) {\t};
  }

  \takeaway{非影像信息既参与构图，也作为独立通道参与编码与融合}
\end{canvas}
\end{frame}

% ============================================================
%  15 致谢
% ============================================================
\begin{frame}[plain]
\begin{canvas}
  \begin{scope}[shift={(7.6,0.2)}]
    \foreach \n/\x/\y in {a/1.0/5.2, b/2.3/6.1, c/3.6/5.4, d/5.2/6.0, e/6.3/4.9,
                          f/0.6/3.6, g/2.0/4.2, h/3.4/3.7, i/4.9/4.2, j/6.0/3.1,
                          k/1.2/1.9, l/2.7/2.4, m/4.1/1.9, n/5.4/1.4, o/3.0/0.8, p/6.3/1.9}
      \coordinate (\n) at (\x,\y);
    \foreach \u/\v/\w/\o in {a/b/1.2/.55, b/c/.6/.35, c/d/1.4/.6, d/e/.7/.4, a/g/.9/.45, b/g/.5/.3,
                             c/h/1.1/.5, c/i/.6/.35, d/i/1.0/.5, e/i/.5/.3, e/j/1.3/.55, f/g/.8/.4,
                             g/h/1.5/.6, h/i/.7/.4, i/j/.9/.45, f/k/1.0/.5, g/l/.6/.35, h/l/1.1/.5,
                             h/m/.6/.3, i/m/.8/.4, j/n/1.0/.5, j/p/.6/.35, k/l/.9/.45, l/m/1.3/.55,
                             m/n/.7/.4, l/o/.6/.35, m/o/1.0/.5, n/p/1.1/.5, k/o/.5/.3}
      \draw[Navy,opacity=\o,line width=\w pt] (\u) -- (\v);
    \foreach \n in {a,c,h,j,l,o,g} \node[case,minimum size=.42cm] at (\n) {};
    \foreach \n in {b,e,f,i,m,p,d,k,n} \node[ctrl,minimum size=.42cm] at (\n) {};
    \draw[Img,line width=1.1pt] (g) circle (.34);
    \draw[Non,line width=1.1pt] (g) circle (.46);
  \end{scope}
  \node[anchor=north west,inner sep=0,font=\fontsize{32}{36}\selectfont\bfseries,text=Navy] at (-.03,5.9) {谢谢聆听};
  \node[anchor=north west,inner sep=0,font=\fontsize{13}{16}\selectfont,text=Muted] at (0,4.25) {欢迎各位老师和同学指正};
  \draw[rule] (0,3.3) -- (6.4,3.3);
  \node[anchor=north west,inner sep=0,align=left,font=\fsmall,text=Ink] at (0,2.75)
    {Cai et al.\ MM-GTUNets: Unified Multi-Modal\\Graph Deep Learning for Brain Disorders Prediction.\\\textit{IEEE Transactions on Medical Imaging}, 2025.};
  \node[anchor=west,inner sep=0,font=\fsmall] at (0,1.05)
    {\textcolor{Muted}{论文}\quad\href{https://arxiv.org/abs/2406.14455}{arxiv.org/abs/2406.14455}};
  \node[anchor=west,inner sep=0,font=\fsmall] at (0,.5)
    {\textcolor{Muted}{代码}\quad\href{https://github.com/NZWANG/MM-GTUNets}{github.com/NZWANG/MM-GTUNets}};
\end{canvas}
\begin{tikzpicture}[remember picture,overlay]
  \node[anchor=north east,inner sep=0] at ([xshift=-.9cm,yshift=-.55cm]current page.north east)
    {\includegraphics[height=1.05cm]{assets/nwpu-crest.png}};
\end{tikzpicture}
\end{frame}
